\section{总结与延伸}

Week 1 建立了整门课程的认知框架：AI 正在从根本上改变软件开发流程，但优秀的工程师需要理解 LLM 的工作原理和局限性，掌握 prompting 技术体系，才能在 human-agent 协作中发挥最大价值。

\subsection{关键要点}

\begin{itemize}
    \item LLM 是基于 Transformer 的 next-token prediction 模型，经过预训练、SFT 和偏好对齐三阶段训练
    \item Prompting 技术包括 zero-shot、k-shot、CoT、self-consistency、tool use、RAG 和 reflexion
    \item 好的 context engineering 是获得好结果的前提
    \item AI 不会取代开发者，但会取代不会使用 AI 的开发者
\end{itemize}

\subsection{拓展阅读}

\begin{itemize}
    \item 课程官网：\url{https://themodernsoftware.dev}
    \item 作业仓库：\url{https://github.com/mihail911/modern-software-dev-assignments}
    \item Anthropic Prompt Engineering Guide：\url{https://docs.anthropic.com/en/docs/build-with-claude/prompt-engineering/prompt-improver}
    \item Vaswani et al., ``Attention Is All You Need'' (2017)
\end{itemize}

%% --- Fin del contenido --- %%

\end{document}
